\subsection{Full Results Table}

Table~\ref{tab:full} lists every DR$\times$clustering combination with four quality metrics
and timing information.

\begin{longtable}{llrrrrrrr}
\toprule
DR & Clustering & DR\,(s) & Clust\,(s) & $K$ & CH & DB & Sil & ARI \\
\midrule
\endfirsthead
\toprule
DR & Clustering & DR\,(s) & Clust\,(s) & $K$ & CH & DB & Sil & ARI \\
\midrule
\endhead
\bottomrule
\endfoot
Raw       & MiniBatchKMeans         &    0.00 &  0.86 & 10 & 1199.5  & 2.122 & 0.146 & 0.343 \\
Raw       & DBSCAN                  &    0.00 &  0.11 &  2 &    6.0  & 1.193 & 0.065 & 0.000 \\
Raw       & AgglomerativeClustering &    0.00 &  3.24 & 10 & 1525.1  & 1.666 & 0.165 & 0.369 \\
Raw       & GaussianMixture         &    0.00 &  6.35 & 10 & 1070.2  & 2.526 & 0.129 & 0.396 \\
Raw       & Birch                   &    0.00 &  2.90 & 10 & 1580.5  & 1.562 & 0.190 & 0.371 \\
\midrule
PCA       & MiniBatchKMeans         &    0.00 &  0.06 & 10 & 1638.0  & 1.774 & 0.185 & 0.318 \\
PCA       & DBSCAN                  &    0.00 &  0.10 &  3 &    4.9  & 1.061 & 0.063 & 0.000 \\
PCA       & AgglomerativeClustering &    0.00 &  3.04 & 10 & 1574.6  & 1.551 & 0.185 & 0.393 \\
PCA       & GaussianMixture         &    0.00 &  5.64 & 10 & 1094.1  & 2.662 & 0.128 & 0.417 \\
PCA       & Birch                   &    0.00 &  2.93 & 10 & 1592.6  & 1.576 & 0.189 & 0.365 \\
\midrule
SAE       & MiniBatchKMeans         &   53.77 &  0.06 & 10 & 1494.7  & 2.076 & 0.127 & 0.348 \\
SAE       & DBSCAN                  &   53.77 &  0.11 &  5 &    4.2  & 1.549 & 0.004 & 0.000 \\
SAE       & AgglomerativeClustering &   53.77 &  3.47 & 10 & 1499.3  & 2.063 & 0.121 & 0.359 \\
SAE       & GaussianMixture         &   53.77 &  7.82 & 10 & 1431.6  & 2.320 & 0.108 & 0.411 \\
SAE       & Birch                   &   53.77 &  3.33 & 10 & 1499.4  & 2.062 & 0.121 & 0.359 \\
\midrule
CNN-SAE   & MiniBatchKMeans         &  350.48 &  0.06 & 10 &  712.2  & 1.951 & 0.148 & 0.362 \\
CNN-SAE   & DBSCAN                  &  350.48 &  0.12 &  2 & ---     & ---   & ---   & ---   \\
CNN-SAE   & AgglomerativeClustering &  350.48 &  3.23 & 10 &  620.2  & 2.108 & 0.119 & 0.383 \\
CNN-SAE   & GaussianMixture         &  350.48 &  4.06 & 10 &  543.9  & 3.119 & 0.126 & 0.463 \\
CNN-SAE   & Birch                   &  350.48 &  3.19 & 10 &  620.2  & 2.108 & 0.119 & 0.383 \\
\midrule
t-SNE     & MiniBatchKMeans         &   97.20 &  0.04 & 10 & 12016.3 & 0.811 & 0.421 & 0.441 \\
t-SNE     & DBSCAN                  &   97.20 &  0.03 & 115 & 4126.9 & 0.877 & 0.034 & 0.260 \\
t-SNE     & AgglomerativeClustering &   97.20 &  2.02 & 10 & 11278.6 & 0.746 & 0.424 & 0.401 \\
t-SNE     & GaussianMixture         &   97.20 &  0.12 & 10 & 10667.8 & 0.839 & 0.403 & 0.434 \\
t-SNE     & Birch                   &   97.20 &  0.40 & 10 & 11609.9 & 0.765 & 0.422 & 0.420 \\
\midrule
UMAP      & MiniBatchKMeans         &   97.33 &  0.06 & 10 & 33793.2 & 0.702 & \textbf{0.541} & \textbf{0.494} \\
UMAP      & DBSCAN                  &   97.33 &  0.09 & 26 &  5334.2 & 1.227 & -0.087 & 0.387 \\
UMAP      & AgglomerativeClustering &   97.33 &  2.61 & 10 & 35379.7 & 0.740 & 0.508 & 0.459 \\
UMAP      & GaussianMixture         &   97.33 &  2.12 & 10 & 29864.0 & 0.806 & 0.462 & 0.421 \\
UMAP      & Birch                   &   97.33 &  0.16 & 10 & \textbf{35898.9} & \textbf{0.658} & 0.540 & 0.463 \\
\bottomrule
\caption{Full results. Bold values mark the best result per metric across all combinations.
  $K$ = number of clusters found. CH = Calinski--Harabasz ($\uparrow$),
  DB = Davies--Bouldin ($\downarrow$), Sil = Silhouette ($\uparrow$), ARI = Adjusted Rand ($\uparrow$).
  t-SNE DR time reflects the original fit (97\,s); subsequent runs loaded from cache.}
\label{tab:full}
\end{longtable}

\subsection{Best Configurations per Metric}

\begin{table}[H]
  \centering
  \begin{tabular}{llll}
    \toprule
    Metric & DR & Clustering & Value \\
    \midrule
    Silhouette ($\uparrow$)           & UMAP & MiniBatchKMeans & 0.541 \\
    Calinski--Harabasz ($\uparrow$)   & UMAP & Birch           & 35\,899 \\
    Davies--Bouldin ($\downarrow$)    & UMAP & Birch           & 0.658 \\
    Adjusted Rand ($\uparrow$)        & UMAP & MiniBatchKMeans & 0.494 \\
    \bottomrule
  \end{tabular}
  \caption{Single-metric winners. UMAP dominates all four metrics.}
  \label{tab:best}
\end{table}

\subsection{Heatmaps}

\begin{figure}[H]
  \centering
  \includegraphics[width=\linewidth]{heatmap_silhouette.png}
  \caption{Silhouette score heatmap (DR rows $\times$ clustering columns). Every algorithm
    improves dramatically when fed a t-SNE or UMAP embedding. DBSCAN is the only algorithm
    that fails even on UMAP, because it keeps labelling most points as noise.}
  \label{fig:heatmap_sil}
\end{figure}

\subsection{Timing Analysis}

DR training time dominates total cost. CNN-SAE is the heaviest ($\approx$6 min on GPU),
followed by t-SNE and UMAP ($\approx$1.6 min each) and SAE ($\approx$54\,s); PCA is
essentially free. Clustering itself is always cheap by comparison (under 10\,s per run
for the slowest combinations).

A second win from non-linear DR: because t-SNE and UMAP produce compact low-dimensional
embeddings, every downstream clustering algorithm runs noticeably faster than on Raw
(784-D) or SAE (64-D) features. So the expensive DR step pays for itself twice —
better metrics \emph{and} faster clustering.

\subsection{Heatmap Observations}

The metric heatmaps tell a consistent story: UMAP and t-SNE dominate every internal metric
by a wide margin, because their non-linear embeddings push same-class points into tight,
well-separated groups. That geometry directly inflates Silhouette and Calinski--Harabasz and
shrinks Davies--Bouldin. Linear and reconstruction-based methods (Raw, PCA, SAE, CNN-SAE) live
in much higher-dimensional spaces where distances are nearly uniform (the curse of
dimensionality), so all internal metrics look mediocre even when the clustering is reasonable.

On Adjusted Rand --- the only metric that compares against the true Fashion-MNIST labels ---
UMAP + MiniBatchKMeans wins outright (ARI $\approx$ 0.494), with CNN-SAE + GMM remaining
competitive (ARI $\approx$ 0.463), suggesting that convolutional features genuinely capture
class structure even without explicit supervision. DBSCAN is the clear loser across the board:
with auto-tuned $\varepsilon$ it either collapses most points into a single cluster or
over-segments into noise, giving near-zero Silhouette and ARI everywhere except t-SNE / UMAP.
